\begin{thebibliography}{99}\small\setlength{\baselineskip}{11pt}\setlength{\itemsep}{0pt}\setlength{\parsep}{0pt}

\bibitem{kagey131}
Kagey, P.: Problem 131. Open Problem Collection.
\url{https://peterkagey.com/problems/131/}. Accessed October 2026

\bibitem{paperA}
Pemmasani, A.: Uniform Bessel bounds for magnetization crossings in finite
Ising chains. Preprint at \url{https://arxiv.org/abs/2610.06935} (2026)

\bibitem{paperB}
Pemmasani, A.: Vertex crossings in a symmetric Markov multinomial model.
Preprint at \url{https://arxiv.org/abs/2610.09258} (2026)

\bibitem{godrecheluck2001}
Godr\`eche, C., Luck, J.M.: Statistics of the occupation time of renewal
processes. J. Stat. Phys. \textbf{104}(3--4), 489--524 (2001).
\url{https://doi.org/10.1023/A:1010364003250}. Section and equation numbers
refer to arXiv:cond-mat/0010428v1

\bibitem{dlmf}
Olver, F.W.J., et al. (eds.): NIST Digital Library of Mathematical Functions.
Section~5.4(iii) and Eqs.~5.4.12, 5.4.14, 5.6.4 and~5.15.5.
\url{https://dlmf.nist.gov}

\bibitem{ds2007}
Dietz, Z., Sethuraman, S.: Occupation laws for some time-nonhomogeneous
Markov chains. Electron. J. Probab. \textbf{12}, 661--683 (2007).
\url{https://doi.org/10.1214/EJP.v12-413}. Example and theorem numbers refer
to arXiv:math/0701798v1

\bibitem{ev2018}
Engl\"ander, J., Volkov, S.: Turning a coin over instead of tossing it. J.
Theor. Probab. \textbf{31}(2), 1097--1118 (2018).
\url{https://doi.org/10.1007/s10959-016-0725-1}. Theorem numbers refer to
arXiv:1606.03281v1

\bibitem{abt2003}
Arratia, R., Barbour, A.D., Tavar\'e, S.: Logarithmic Combinatorial
Structures: A Probabilistic Approach. EMS Monographs in Mathematics. European
Mathematical Society, Z\"urich (2003). \url{https://doi.org/10.4171/000}

\bibitem{bc2018}
Bouguet, F., Cloez, B.: Fluctuations of the empirical measure of freezing
Markov chains. Electron. J. Probab. \textbf{23}, 2 (2018).
\url{https://doi.org/10.1214/17-EJP130}. Proposition numbers refer to
arXiv:1705.02121v1

\bibitem{dls2023}
Dietz, Z., Lippitt, W., Sethuraman, S.: Stick-breaking processes, clumping,
and Markov chain occupation laws. Sankhy\=a A \textbf{85}(1), 129--171
(2023). \url{https://doi.org/10.1007/s13171-020-00236-x}. Theorem numbers
refer to arXiv:1901.08135v1

\bibitem{ev2022}
Engl\"ander, J., Volkov, S.: Conservative random walk. Electron. J. Probab.
\textbf{27}, 138 (2022). \url{https://doi.org/10.1214/22-EJP863}

\bibitem{ev2024}
Engl\"ander, J., Volkov, S.: Coin-Turning, Random Walks and Inhomogeneous
Markov Chains. Author draft of March 18, 2024.
\url{https://www.researchgate.net/publication/379333206_Coin-Turning_Random_Walks_and_Inhomogeneous_Markov_Chains}.
Chapter, section, theorem and page numbers in this note refer to this draft.
We did not check them against the published edition~\cite{ev2025}

\bibitem{ev2025}
Engl\"ander, J., Volkov, S.: Coin-Turning, Random Walks and Inhomogeneous
Markov Chains. World Scientific, Singapore (2025).
\url{https://doi.org/10.1142/13837}

\bibitem{evw2020}
Engl\"ander, J., Volkov, S., Wang, Z.: The coin-turning walk and its scaling
limit. Electron. J. Probab. \textbf{25}, 3 (2020).
\url{https://doi.org/10.1214/19-EJP406}

\bibitem{hoeffding1963}
Hoeffding, W.: Probability inequalities for sums of bounded random variables.
J. Am. Stat. Assoc. \textbf{58}(301), 13--30 (1963).
\url{https://doi.org/10.1080/01621459.1963.10500830}

\bibitem{prr2026}
Pradhan, A., Roy, R., Ray, P.: Diffusive-to-ballistic transition in a
persistent random walk. Preprint at \url{https://arxiv.org/abs/2605.20116}
(2026)

\end{thebibliography}
